\subsection{Abordagem Educomunicativa}

\subsubsection{Abordagem Educomunicativa}

\paragraph{Fundamentação Teórica}

\subparagraph{O que é Educomunicação?}

A \textbf{educomunicação} é um campo interdisciplinar que integra comunicação e educação, reconhecendo que:

\begin{quote}
\textit{``A comunicação é um processo educativo e a educação é um processo comunicativo''} --- Soares e Melo (2020)
\end{quote}

No Projeto Hermathena, adotamos a educomunicação como eixo central porque ela permite:

\begin{itemize}
\item \textbf{Diálogo horizontal} entre cientistas e público
\item \textbf{Construção conjunta} de conhecimento
\item \textbf{Contextualização} do conteúdo científico
\item \textbf{Empoderamento} dos participantes
\end{itemize}

\subparagraph{Três Vertentes Teóricas}

\begin{verbatim}
┌─────────────────────────────────────────────────────────────┐
│                 FUNDAÇÃO TEÓRICA                           │
├─────────────────────────────────────────────────────────────┤
│                                                             │
│   ┌──────────────┐  ┌──────────────┐  ┌──────────────┐    │
│   │ Educomuni -   │  │ Comunicação  │  │ Desenvolvi -  │    │
│   │ cação        │  │ Pública da   │  │ mento        │    │
│   │              │  │ Ciência      │  │ Regional     │    │
│   └──────┬───────┘  └──────┬───────┘  └──────┬───────┘    │
│          │                 │                 │             │
│          └─────────────────┼─────────────────┘             │
│                            │                               │
│                            ▼                               │
│                 ┌──────────────────────┐                   │
│                 │  PROJETO HERMATHENA  │                   │
│                 └──────────────────────┘                   │
│                                                             │
└─────────────────────────────────────────────────────────────┘
\end{verbatim}


\bigskip
\centerline{\rule{13cm}{0.4pt}}
\bigskip

\paragraph{1. Educomunicação}

\subparagraph{Princípios Fundamentais}

\subparagraph{1.1 Diálogo Horizontal}

Contrário ao modelo tradicional de transmissão de informação, propomos:

\bigskip\noindent
\begin{tabular}{p{\dimexpr 0.500\linewidth-2\tabcolsep}p{\dimexpr 0.500\linewidth-2\tabcolsep}}
\toprule
Modelo Tradicional & Modelo Educomunicativo \\
\hline
Cientista ensina & Cientista e público aprendem juntos \\
Conhecimento unidirecional & Troca de saberes \\
Público passivo & Público ativo e produtor \\
Linguagem técnica & Linguagem contextualizada \\
\bottomrule
\end{tabular}

\bigskip

\subparagraph{1.2 Processo Dialógico}

\begin{verbatim}
┌──────────────────────────────────────────────────────────────┐
│                     CICLO DIALÓGICO                         │
├──────────────────────────────────────────────────────────────┤
│                                                              │
│   ┌──────────┐      ┌──────────┐      ┌──────────┐         │
│   │  OUÇA    │ ───> │  DIALOGUE│ ───> │  CO -RIE  │         │
│   │          │      │          │      │          │         │
│   │ Entenda  │      │ Troque   │      │ Construa │         │
│   │ necessi - │      │ saberes  │      │ junto    │         │
│   │ dades    │      │          │      │          │         │
│   └──────────┘      └──────────┘      └────┬─────┘         │
│         ^                                   │               │
│         │                                   │               │
│         └───────────────────────────────────┘               │
│                        AVALIE                               │
│                                                              │
└──────────────────────────────────────────────────────────────┘
\end{verbatim}

\subparagraph{Aplicação no Projeto}

\textbf{Exemplo: Produção de Material Didático}

\begin{enumerate}
\item \textbf{Pesquisa participativa}

\begin{itemize}
\item Entrevistas com professores sobre necessidades
\item Oficinas com estudantes para entender interesses
\item Consulta a especialistas em educação
\end{itemize}


\item \textbf{Co-criação}

\begin{itemize}
\item Workshops de roteirização com educadores
\item Testes pilotos em sala de aula
\item Feedback contínuo
\end{itemize}


\item \textbf{Validação comunitária}

\begin{itemize}
\item Revisão por professores da rede pública
\item Ajustes baseados na realidade escolar
\item Aprovação por equipe científica
\end{itemize}
\end{enumerate}


\bigskip
\centerline{\rule{13cm}{0.4pt}}
\bigskip

\paragraph{2. Comunicação Pública da Ciência}

\subparagraph{Modelo de Diálogo entre Iguais}

Segundo Irwin (2008), a comunicação pública da ciência evoluiu:

\bigskip\noindent
\begin{tabular}{p{\dimexpr 0.333\linewidth-2\tabcolsep}p{\dimexpr 0.333\linewidth-2\tabcolsep}p{\dimexpr 0.333\linewidth-2\tabcolsep}}
\toprule
Modelo & Característica & Era \\
\hline
\textbf{Deficitário} & ``O público precisa ser educado'' & 1950-1980 \\
\textbf{Diálogo} & ``Troca de perspectivas'' & 1990-2010 \\
\textbf{Participativo} & ``Ciência com e para a sociedade'' & 2010-presente \\
\bottomrule
\end{tabular}

\bigskip

\subparagraph{Abordagem do Hermathena: Modelo Participativo}

Características:

\begin{itemize}
\item \textbf{Reconhecimento de expertise local}: A comunidade tem saberes legítimos
\item \textbf{Transparência científica}: Divulgação de incertezas e limitações
\item \textbf{Deliberação pública}: Espaço para questionamentos
\item \textbf{Responsividade}: Ciência responde às demandas sociais
\end{itemize}

\subparagraph{Estratégias de Implementação}

\subparagraph{2.1 Ciência em Contexto}

Conectar o conteúdo à realidade do Sertão paraibano:

\bigskip\noindent
\begin{tabular}{p{\dimexpr 0.500\linewidth-2\tabcolsep}p{\dimexpr 0.500\linewidth-2\tabcolsep}}
\toprule
Tema Científico & Conexão Local \\
\hline
Energia escura & A ``força invisível'' como metáfora para o desenvolvimento regional \\
Radioastronomia & Tecnologia aplicada à agricultura (sensores) \\
BAO e escala cósmica & A escala do Sertão para o Brasil e mundo \\
Espectro eletromagnética & Comparação com comunicação por rádio local \\
\bottomrule
\end{tabular}

\bigskip

\subparagraph{2.2 Narrativas Inclusivas}

\begin{itemize}
\item Histórias de paraibanos na ciência
\item Reconhecimento de saberes tradicionais
\item Diálogo com cultura local (cordel, forró, arte sertaneja)
\end{itemize}


\bigskip
\centerline{\rule{13cm}{0.4pt}}
\bigskip

\paragraph{3. Desenvolvimento Regional}

\subparagraph{Ciência como Indutora de Desenvolvimento}

Segundo Lastres et al. (2005), a ciência e tecnologia funcionam como indutores de desenvolvimento quando:

\begin{itemize}
\item Articuladas com políticas públicas
\item Integradas à iniciativa privada
\item Apropriadas pela comunidade local
\end{itemize}

\subparagraph{O BINGO como Vetor de Desenvolvimento}

\begin{verbatim}
┌─────────────────────────────────────────────────────────────┐
│           IMPACTOS DO BINGO NA REGIÃO                       │
├─────────────────────────────────────────────────────────────┤
│                                                             │
│   ┌──────────────┐  ┌──────────────┐  ┌──────────────┐    │
│   │ HUMANO       │  │ CIENTÍFICO   │  │ ECONÔMICO    │    │
│   │              │  │              │  │              │    │
│   │ \cdot Formação   │  │ \cdot Pesquisa   │  │ \cdot Turismo    │    │
│   │   STEM       │  │   de ponta   │  │   científico │    │
│   │ \cdot Empoderam. │  │ \cdot Inovação   │  │ \cdot Inovação   │    │
│   │   comunitário│  │   tecnológica│  │   tecnológica│    │
│   └──────────────┘  └──────────────┘  └──────────────┘    │
│                                                             │
└─────────────────────────────────────────────────────────────┘
\end{verbatim}

\subparagraph{Comunicação como Facilitadora}

A comunicação científica do Hermathena potencializa esses impactos:

\begin{enumerate}
\item \textbf{Visibilidade nacional}: Atrai atenção e recursos para a região
\item \textbf{Educação de qualidade}: Prepara jovens para oportunidades
\item \textbf{Turismo científico}: Gera renda local
\item \textbf{Atração de investimentos}: Conecta empresas à ciência
\end{enumerate}


\bigskip
\centerline{\rule{13cm}{0.4pt}}
\bigskip

\paragraph{Metodologias Específicas}

\subparagraph{1. Design Instrucional}

\subparagraph{Análise de Necessidades (ADDIE)}

\begin{verbatim}
┌─────────┐   ┌─────────┐   ┌─────────┐   ┌─────────┐   ┌─────────┐
│ ANALISAR│ \rightarrow │  DESIGN │ \rightarrow │ DESENV. │ \rightarrow │IMPLEMENT│ \rightarrow │ AVALIAR │
└─────────┘   └─────────┘   └─────────┘   └─────────┘   └─────────┘
    │             │             │             │             │
    ▼             ▼             ▼             ▼             ▼
Perfil do    Objetivos    Produção     Aplicação   Métricas
devedor      de            do           em campo    de sucesso
             aprendizado   material
\end{verbatim}

\subparagraph{Aplicação}

\bigskip\noindent
\begin{tabular}{p{\dimexpr 0.500\linewidth-2\tabcolsep}p{\dimexpr 0.500\linewidth-2\tabcolsep}}
\toprule
Etapa & Ação no Hermathena \\
\hline
\textbf{Análise} & Pesquisa com 50 professores da rede pública \\
\textbf{Design} & Definição de objetivos por faixa etária \\
\textbf{Desenvolvimento} & Criação de materiais com equipe multidisciplinar \\
\textbf{Implementação} & Teste em 10 escolas piloto \\
\textbf{Avaliação} & Medição de aprendizado e satisfação \\
\bottomrule
\end{tabular}

\bigskip

\subparagraph{2. Produção Colaborativa}

\subparagraph{Modelo de Produção em Rede}

\begin{verbatim}
┌───────────────────┐
                    │   COORDENAÇÃO     │
                    │   HERMATHENA      │
                    └─────────┬─────────┘
                              │
        ┌─────────────────────┼─────────────────────┐
        │                     │                     │
        ▼                     ▼                     ▼
┌───────────────┐    ┌───────────────┐    ┌───────────────┐
│   ESCOLAS     │    │  UNIVERSIDADES│    │  COMUNIDADE   │
│               │    │               │    │               │
│ \cdot Professores │    │ \cdot Pesquisadores│   │ \cdot Lideranças  │
│ \cdot Estudantes  │    │ \cdot Graduandos  │    │ \cdot Artistas    │
│ \cdot Gestores    │    │ \cdot Pós -grad.   │    │ \cdot Empreended. │
└───────────────┘    └───────────────┘    └───────────────┘
\end{verbatim}

\subparagraph{3. Avaliação Participativa}

\subparagraph{Instrumentos}

\begin{enumerate}
\item \textbf{Pesquisa de Percepção Comunitária}

\begin{itemize}
\item Aplicação bianual
\item Amostra: 500 habitantes de Aguiar e região
\item Foco: Reconhecimento do BINGO, interesse em ciência
\end{itemize}


\item \textbf{Questionário de Satisfação}

\begin{itemize}
\item Pós-cada atividade
\item Escala Likert (1-5)
\item Espaço para sugestões abertas
\end{itemize}


\item \textbf{Grupos Focais}

\begin{itemize}
\item 2x por ano
\item 8-12 participantes
\item Discussão guiada sobre impacto
\end{itemize}


\item \textbf{Estudos de Caso}

\begin{itemize}
\item Narrativas individuais
\item Histórias de transformação
\item Documentação qualitativa
\end{itemize}
\end{enumerate}


\bigskip
\centerline{\rule{13cm}{0.4pt}}
\bigskip

\paragraph{Matriz de Conteúdos por Público}

\subparagraph{Ensino Fundamental (Anos Finais)}

\bigskip\noindent
\begin{tabular}{p{\dimexpr 0.333\linewidth-2\tabcolsep}p{\dimexpr 0.333\linewidth-2\tabcolsep}p{\dimexpr 0.333\linewidth-2\tabcolsep}}
\toprule
Tema & Abordagem & Produto \\
\hline
Sistema Solar & Jogos e atividades lúdicas & Cartilha ilustrada \\
Estrelas & Contação de histórias & Vídeos animados \\
Universo & Maquetes e modelos & Kit escolar \\
\bottomrule
\end{tabular}

\bigskip

\subparagraph{Ensino Médio}

\bigskip\noindent
\begin{tabular}{p{\dimexpr 0.333\linewidth-2\tabcolsep}p{\dimexpr 0.333\linewidth-2\tabcolsep}p{\dimexpr 0.333\linewidth-2\tabcolsep}}
\toprule
Tema & Abordagem & Produto \\
\hline
Cosmologia & Simulações computacionais & Software educativo \\
Radioastronomia & Experimentos práticos & Manual laboratorial \\
Método científico & Projetos de pesquisa & Guia do estudante \\
\bottomrule
\end{tabular}

\bigskip

\subparagraph{Educadores}

\bigskip\noindent
\begin{tabular}{p{\dimexpr 0.333\linewidth-2\tabcolsep}p{\dimexpr 0.333\linewidth-2\tabcolsep}p{\dimexpr 0.333\linewidth-2\tabcolsep}}
\toprule
Tema & Abordagem & Produto \\
\hline
Astronomia na sala & Formação continuada & Curso online \\
Uso de materiais didáticos & Oficinas práticas & Tutoriais em vídeo \\
Avaliação em ciências & Estudos de caso & Caderno de práticas \\
\bottomrule
\end{tabular}

\bigskip

\subparagraph{Público Geral}

\bigskip\noindent
\begin{tabular}{p{\dimexpr 0.333\linewidth-2\tabcolsep}p{\dimexpr 0.333\linewidth-2\tabcolsep}p{\dimexpr 0.333\linewidth-2\tabcolsep}}
\toprule
Tema & Abordagem & Produto \\
\hline
BINGO e a Paraíba & Documentários & Vídeos no YouTube \\
Curiosidades astronômicas & Redes sociais & Posts e reels \\
Ciência no cotidiano & Podcasts & Episódios mensais \\
\bottomrule
\end{tabular}

\bigskip


\bigskip
\centerline{\rule{13cm}{0.4pt}}
\bigskip

\paragraph{Cronograma Metodológico}

\subparagraph{Fase 1: Diagnóstico e Planejamento (Meses 1-2)}

\bigskip\noindent
\begin{tabular}{p{\dimexpr 0.333\linewidth-2\tabcolsep}p{\dimexpr 0.333\linewidth-2\tabcolsep}p{\dimexpr 0.333\linewidth-2\tabcolsep}}
\toprule
Semana & Atividade & Metodologia \\
\hline
1-2 & Pesquisa de campo com escolas & Entrevistas, questionários \\
3-4 & Análise de necessidades & Grupos focais \\
5-6 & Definição de matriz de conteúdos & Workshop participativo \\
7-8 & Planejamento pedagógico & Design instrucional \\
\bottomrule
\end{tabular}

\bigskip

\subparagraph{Fase 2: Produção e Testes (Meses 3-6)}

\bigskip\noindent
\begin{tabular}{p{\dimexpr 0.333\linewidth-2\tabcolsep}p{\dimexpr 0.333\linewidth-2\tabcolsep}p{\dimexpr 0.333\linewidth-2\tabcolsep}}
\toprule
Semana & Atividade & Metodologia \\
\hline
9-12 & Produção de materiais v1 & Co-criação com educadores \\
13-16 & Testes pilotos & Aplicação em 10 escolas \\
17-20 & Ajustes e refinamento & Análise de feedback \\
21-24 & Validação final & Comitê de especialistas \\
\bottomrule
\end{tabular}

\bigskip

\subparagraph{Fase 3: Implementação e Avaliação (Meses 7-12)}

\bigskip\noindent
\begin{tabular}{p{\dimexpr 0.333\linewidth-2\tabcolsep}p{\dimexpr 0.333\linewidth-2\tabcolsep}p{\dimexpr 0.333\linewidth-2\tabcolsep}}
\toprule
Semana & Atividade & Metodologia \\
\hline
25-32 & Implementação em larga escala & Ações nas escolas \\
33-40 & Monitoramento contínuo & Coleta de métricas \\
41-48 & Avaliação de impacto & Pesquisa de percepção \\
49-52 & Documentação e relatório & Análise de dados \\
\bottomrule
\end{tabular}

\bigskip


\bigskip
\centerline{\rule{13cm}{0.4pt}}
\bigskip

\paragraph{Referências}

\begin{itemize}
\item ALONSO, L. et al. \textit{A Educomunicação nos processos de educação e difusão de conhecimento: divulgação científica no Projeto Telescópio BINGO}. ECA-USP, 2023.
\item IRWIN, A. \textit{Risk, science and public communication}. In: The handbook of science and technology studies. 2008.
\item LASTRES, H. M. M.; CASSIOLATO, J. E.; ARROIO, A. (Org.). \textit{Conhecimento, sistemas de inovação e desenvolvimento}. Editora UFRJ, 2005.
\item SOARES, I. M. de M.; MELO, J. M. de. \textit{Educomunicação: conceitos e práticas}. Editora PUC-Rio, 2020.
\end{itemize}


\bigskip
\centerline{\rule{13cm}{0.4pt}}
\bigskip

\textit{Metodologia adaptável conforme aprendizados do projeto.}